\chapter{Benchmarking Cryptographic Primitives}

\section{Operation-Level Benchmarks}

Operation-level benchmarks isolate primitive kernels inside the zkVM and measure
their contribution to proving cost. The workload set focuses on operations that
are expected to dominate cryptographic traces, including substitution layers,
matrix-heavy transforms, and bitwise round-function kernels.

The key engineering heuristic motivating this benchmark set is that the S-box
is typically the most expensive component to prove in a block cipher, because it
is a nonlinear function that must be encoded as a non-trivial constraint in the
execution trace. If a cipher's S-box is costly to evaluate in the VM, swapping
to a different S-box implementation --- for example, from a lookup-table
representation to a bitsliced logic circuit, or vice versa --- could
significantly change the overall proving time. Operation-level benchmarks make
this cost visible before committing to a full-cipher implementation strategy. In
AES, for instance, the SubBytes step applies a single 8-to-8-bit S-box to every
byte of state in every round; identifying whether the S-box or the linear layer
dominates trace cost helps direct optimization effort to where it will have the
most impact.


This focus is motivated by proof-system cost structure. While RISC Zero is a
VM-centric proving stack rather than a hand-authored arithmetic circuit, prover
work still depends on enforcing constraints over the execution trace. In
practice, operation shape and data movement patterns can increase trace pressure
and therefore proving time \cite{bensasson2018stark,risczeroTraceDocs}.

These measurements are used to rank optimization direction, then cross-checked
against full cipher runs. They are evidence for engineering priority, not
standalone performance claims about complete cryptographic pipelines.

\section{Cipher-Level Benchmarks}

Cipher-level benchmarks execute complete implementations under the same RISC Zero
proving flow. The core comparison is AES, Salsa20, and LowMC baselines, with
optimized variants analyzed separately.

\paragraph{AES-192.}
AES is a standardized block cipher with a 128-bit block size; this project uses
the 192-bit-key variant \cite{nistfips197,daemen2002rijndael}. A block cipher
maps one fixed-size plaintext block to one ciphertext block under a secret key.
AES does this by maintaining a compact 16-byte state and applying repeated
rounds of byte substitution, permutation, linear mixing, and key addition. The
main round steps are SubBytes, which applies a nonlinear 8-bit S-box to each
byte; ShiftRows and MixColumns, which diffuse bytes across the state; and
AddRoundKey, which XORs in key-schedule material. In this dissertation AES is
used first as a complete primitive benchmark and later as the block function
inside CTR mode, where AES encrypts counter blocks to produce a keystream. In
the zkVM, the relevant cost question is not only the abstract cipher design, but
how these byte-level table lookups, XORs, and mixing operations compile into
guest instructions.

\paragraph{Salsa20.}
Salsa20 is a stream cipher built from ARX operations: modular addition, bitwise
rotation, and XOR \cite{bernstein2008salsa20}. Instead of encrypting fixed-size
blocks directly like AES, it generates a keystream from a key, nonce, and counter,
then XORs that keystream with the message. Internally, Salsa20 updates a small
word-oriented state through repeated quarter-round functions, where each
quarter-round mixes four 32-bit words using only addition modulo \(2^{32}\),
rotation, and XOR. This makes the implementation regular and software-friendly:
there are no S-box tables and no large binary matrices, but there are many
repeated word operations. Salsa20 is therefore a useful midpoint in this
dissertation: it is not designed specifically for ZK circuits, but it avoids the
lookup-heavy structure of AES and the matrix-heavy structure of LowMC.

\paragraph{LowMC.}
LowMC is a block cipher designed for settings such as MPC and FHE, where
nonlinear gates are expensive \cite{albrecht2015ciphers}. It deliberately uses a
small nonlinear layer and relies heavily on binary linear layers, especially
matrix--vector products over GF(2). A LowMC round applies a limited number of
small S-boxes, multiplies the full state by a binary matrix, and adds round-key
material. The intended benefit is that the scarce nonlinear part keeps
multiplicative complexity low in hand-authored circuits. The tradeoff is that a
software implementation still has to represent and apply the linear matrices,
usually through packed words, bit operations, loops, and memory accesses. This
gives LowMC a strong circuit-level motivation, but it also creates the main
empirical question for this project: whether matrix-heavy software execution
inside RISC Zero is actually cheaper than the more conventional AES path.

Interpretation rule: zkVM proving cost does not necessarily mirror native CPU
performance intuition, so conclusions are based on measured prover behavior.
Where operation-level and cipher-level conclusions disagree, the cipher-level
result is treated as authoritative for system-design decisions.

\section{Cross-Algorithm Results}

Cross-algorithm comparison values used for the headline primitive choice are
taken from the latest commit-locked main report run. That run includes
\texttt{aes-r0} and \texttt{lowmc-r0-optimised}; older Salsa20 and unoptimised
LowMC snapshots remain useful exploratory context, but they are not treated as
the latest headline evidence.

Salsa20 is therefore not absent from the artifact trail, but it is separated from
the headline table. The older support run
\texttt{artifacts/benchmarks/20260305-145335Z} measured \texttt{salsa-r0} over
three successful trials, with median proving time 30.7318 s, median verification
time 0.0118 s, and median user cycles 467020 for a 1024-byte payload. Because
that run used an older campaign and a different payload shape, these values are
used only to document the implemented Salsa20 baseline, not for the final
AES-versus-LowMC ranking in Table~\ref{tab:cross-algo-comparison}.

\begin{table}[H]
\centering
\reporttablesetup
\begin{tabular*}{\textwidth}{@{\extracolsep{\fill}}lrrrr@{}}
\hline
Benchmark ID & Prove s (median) & Verify s (median) & User cycles (median) & Success \\
\hline
\texttt{aes-r0} & 14.530 & 0.0111 & 219445 & 5/5 \\
\texttt{lowmc-r0-optimised} & 484.279 & 0.0970 & 7589089 & 5/5 \\
\hline
\end{tabular*}
\caption{Headline AES versus optimised LowMC comparison from the commit-locked
\(N=5\) main report run.}
\label{tab:cross-algo-comparison}
\end{table}

The effect-size gap is large: optimised LowMC remains substantially slower than
AES even under the latest five-trial report policy.

The important interpretation is not merely that AES is faster in Table
\ref{tab:cross-algo-comparison}. The result shows that the naive ordering
suggested by circuit-friendliness is not preserved by this VM implementation:
LowMC has the strongest ZK-oriented design motivation, but the latest optimised
target still proves much more slowly than the compact AES baseline in this
specific stack.

This also explains how the result is used later. The table is not a universal
claim about primitive quality; it is a decision input for the implemented RISC
Zero pipeline. It justifies pivoting the final authenticated construction toward
AES while keeping LowMC analysis as evidence for why VM-level measurements are
needed.

\section{Operations Directional Ranking}

This directional ranking uses benchmark IDs
\texttt{operation-bnchmrk-r0:*} from the latest operations campaign. Trial
policy is five trials per operation row, and all listed rows succeeded
(\texttt{success\_rate = 1.0}).

The operation rows separate into coarse trace buckets. Addition and subtraction
fall into the 32768-cycle bucket and prove in about 1.8--1.9 seconds, while the
boolean, rotation, multiplication, and Toffoli-style rows fall into the
65536-cycle bucket and cluster around 3.6--3.7 seconds. Within that larger bucket,
rotations and multiplication have the highest user-cycle counts, so they are
useful hotspot indicators even though the rounded trace size hides some fine
ordering in wall-clock time.

\begin{figure}[t]
\centering
\includegraphics[width=0.95\linewidth]{Figures/report-benchmarks/operations_directional_prove_time.png}
\caption{Operation-level directional ranking by median proving time from the report operations campaign.}
\label{fig:operations-directional-prove}
\end{figure}

\section{Operation-to-Cipher Link Note}

The operation campaign and cipher campaign align on one key expectation:
lightweight, individually cheap kernels are not enough to make a primitive
zkVM-efficient by themselves. Full-cipher composition overhead dominates.

The results in Figure~\ref{fig:operations-directional-prove} show that
addition and subtraction are the cheapest operations, while rotations,
multiplication, and boolean operations are more expensive per invocation. A
cipher built primarily from lightweight operations would, in isolation, be
expected to prove quickly. However, individually cheaper operations do not
guarantee a cheap full-cipher result when they are repeated many times per round
and move substantial state through the guest program.

The key insight is therefore compositional: operation rows can identify plausible
pressure points, and LowMC's matrix-heavy structure is consistent with its high
full-cipher cost, but operation-level evidence is insufficient for final
primitive selection.
